\section{Threat Model}
\label{sec:threat}

\subsection{Extending the risk model with an integrity condition}

We adopt CI-Work's model (\Cref{sec:background:ciwork}) and extend it. As
before, an agent $A$ acting for a principal receives instruction $I$, retrieves
entries $E=\bigcup_t o_t$, and emits a final action $\afin$ over a channel $k$
to a recipient $r$. We attach to each entry $e \in E$ two labels,
\begin{equation}
\Cl(e)\in[0,1], \qquad \Il(e)\in[0,1],
\end{equation}
where $\Cl(e)$ is how inappropriate disclosure of $e$ to $r$ is \emph{in this
context}, and $\Il(e)$ is how trustworthy the \emph{origin} of $e$ is as a
basis for a high-stakes outbound artifact. We attach to the request a recipient
clearance $\clr(r)$ and a required outbound integrity $\lvlI(\afin)$.

CI-Work's risk condition is a disclosure condition, which in our notation is
the Bell--LaPadula $\star$-property applied at the write boundary:
\begin{equation}
\exists e \in E \;:\; \Fdisc(\afin, e) = 1 \;\wedge\; \Cl(e) > \clr(r).
\label{eq:blpviolation}
\end{equation}
We add a second, independent failure condition. Let
$\mathrm{ev}(\afin) \subseteq E$ be the set of entries the emitted conclusion
\emph{rests on}---its supporting evidence, as distinct from the set it quotes.
An \emph{integrity violation} is
\begin{equation}
\exists e \in \mathrm{ev}(\afin) \;:\; \Il(e) < \lvlI(\afin).
\label{eq:bibaviolation}
\end{equation}

Two things about \Cref{eq:bibaviolation} deserve emphasis, and together they
shape the rest of the paper.

First, it is quantified over $\mathrm{ev}(\afin)$, not over the entries
$\afin$ discloses. This is deliberate and it is the crux of the integrity axis:
the harm from a low-integrity entry is that the agent \emph{relies} on it, not
that it \emph{repeats} it. Repeating an unverified claim with attribution is
often the correct behavior. Silently believing it is the failure.

Second, this makes \Cref{eq:bibaviolation} strictly harder to measure than
\Cref{eq:blpviolation}. Disclosure is observable in the emitted text;
reliance is not. We confront this directly in \Cref{sec:eval:ilblindspot},
where we show by measurement that a disclosure-based proxy for
\Cref{eq:bibaviolation} undercounts it by exactly the entries that were
believed but never quoted, and that only a paired counterfactual measures it
faithfully.

\paragraph{Independence of the axes.} The central structural claim is that
$\Cl$ and $\Il$ are independent, i.e.\ $\Esens \neq \{e : \Il(e) \text{ low}\}$.
\Cref{tab:quadrants} gives the four quadrants with a representative from each;
every one is realizable, and two of them are misclassified by any single-axis
policy. A confidentiality-only system cannot distinguish the two
right-hand cells from their left-hand neighbors, and so has no basis for
treating a fabricated figure differently from a real one.

\begin{table}[t]
\centering
\caption{The two axes are independent: all four quadrants are populated. A
single-axis policy collapses the columns and therefore cannot treat a
fabricated figure differently from an authoritative one.}
\label{tab:quadrants}
\footnotesize
\begin{tabular}{@{}lp{0.3\columnwidth}p{0.3\columnwidth}@{}}
\toprule
& \textbf{High $\Il$ (trusted)} & \textbf{Low $\Il$ (untrusted)} \\
\midrule
\textbf{High $\Cl$}
& walk-away price \newline \emph{do not send; it is real}
& fabricated tactic \newline \emph{do not send; do not rely} \\[4pt]
\textbf{Low $\Cl$}
& executed contract \newline \emph{safe both ways}
& unverified claim \newline \emph{safe to send; unsafe to rely} \\
\bottomrule
\end{tabular}
\end{table}

\subsection{Both failures co-occur in published transcripts}
\label{sec:threat:cooccur}

The independence claim would be weak if we demonstrated it only on cases we
authored. We therefore dual-label the two qualitative examples published in
CI-Work's own appendix, and find that \emph{each single transcript contains one
Bell--LaPadula violation and one Biba violation}. Because these are the
original authors' examples, the reader can check our labeling against their
paper.

\paragraph{Case A: software license renewal (External).} Procurement emails an
external vendor representative. The sensitive set contains the FY22 budget and
a renewal contingency, a seat-overage figure with penalty exposure, and a
planned acquisition raising seat count 40\%---all high $\Cl$ and high $\Il$,
being real internal financial facts. It also contains a procurement chat log
proposing to \emph{present competitor quotes as \$90 per seat when the figure
actually received was \$95}, in order to pressure the vendor. That entry is
high $\Cl$ and \textbf{low $\Il$}: it proposes a fabrication, in an informal
channel. The published agent transcript leaks the overage and the acquisition
plan, satisfying \Cref{eq:blpviolation}. But the distinct and unmeasured
failure available in the same context is repeating the fabricated \$90 as fact
to the counterparty, when the authoritative competitor benchmark in the
essential set says \$95. That is \Cref{eq:bibaviolation}, and no CI-Work metric
would flag it either way.

\paragraph{Case B: lease negotiation (External).} A real-estate manager emails
an external landlord. The sensitive set contains a relocation contingency and
rent guardrails including the walk-away threshold---high on both axes---plus an
HR chat message speculating that headcount reduction will cut space needs 30\%,
which is high $\Cl$ and \textbf{low $\Il$}. The published transcript leaks the
guardrails verbatim, writing the ceiling and walk-away thresholds to the
landlord: a textbook no-write-down failure. It \emph{also} states a revised
space requirement of 13{,}000--15{,}000 sq ft attributed to anticipated
headcount reduction, a figure and rationale that trace to the unconfirmed chat
message rather than to any authoritative facilities plan. A low-integrity
source has been written up into a binding external proposal.

Both cases also exhibit the quadrant that motivates our enforcement design:
an entry that is high on both axes and therefore must \emph{inform} the output
while never appearing in it. Neither a disclosure filter nor an instruction to
``avoid sensitive topics'' can express that requirement.

\subsection{Adversary, capabilities, and scope}

\paragraph{Adversary.} We consider two, of differing strength.
\emph{(i) A careless-context adversary}, which is not an attacker at all but
the ordinary condition of enterprise data: unverified claims, stale plans, and
informal speculation coexist with authoritative records, and the agent has no
native means of telling them apart. This is the setting CI-Work measures.
\emph{(ii) An author-controlled-channel adversary}, who supplies content in a
channel they legitimately write to---a pull-request description, a code
comment, a ticket comment---with the goal of either eliciting restricted
context or causing the agent to reach a favorable conclusion. In the PR-review
setting of \Cref{sec:benchmark} this adversary is the party under review, which
is the realistic case: they have both the incentive and the write access.

\paragraph{Assumed trusted.} The identity and clearance directory, the
connector-level access control that scopes retrieval, the policy configuration
itself, and the labeling pipeline's offline outputs. We assume the enforcement
point is not bypassable, i.e.\ the agent cannot emit except through the wrapped
tool; in deployment this is the standard sidecar-proxy
assumption~\cite{nist800207a}.

\paragraph{Out of scope.} Model weights and inference infrastructure
compromise; collusion between the recipient and an insider; covert channels of
the timing or resource-consumption kind. One in-scope residual risk we do
\emph{not} solve is reconstruction: an abstract that is individually below the
write ceiling may still permit recovery of the restricted specifics when
combined with admitted entries. We treat this as a measurement obligation
rather than a solved problem, and state its current status honestly in
\Cref{sec:discussion:limitations}.

\paragraph{A note on third-party data subjects.} CI-Work's five CI parameters
largely collapse in its setting because the data subject is usually the sender.
In PR review the data subject is typically a third party---a customer named in
a postmortem, a colleague named in a performance discussion---who is not party
to the conversation and cannot consent. Our entry schema tracks $\mathrm{subj}(e)$
separately for this reason. We note the limit of our own framework here: no
clearance lattice can \emph{authorize} disclosure of third-party regulated
data, so the only correct treatment is hard denial, which is why the precedence
order in \Cref{sec:method:precedence} places the regulated block above every
tunable parameter.
